\documentstyle[% 


righttag% 


,11pt% 


,amssymb% 


,russian%               remove in Eng. 


,izvru%                 Set izven% in Eng. 


]{amsart} 


 


%  


\newcommand{\const}{\operatorname{const}} 


\newcommand{\tg}{\operatorname{tg}} 


\newcommand{\ctg}{\operatorname{ctg}} 


\newcommand{\tts}[1]{} 


 


\theoremstyle{plain} 


\theorembodyfont{\it} 


\newtheorem{thmnonum}{\hskip\parindent Теорема} 


\renewcommand{\thethmnonum}{} 


\renewcommand{\baselinestretch}{0.96}


 


%\hoffset=-15mm%                  For Eng. 


 


\setcounter{page}{37} 


\god{2000} 


\nomer{7~(458)} %                        Remove in Eng. 


%\nomer{Vol.\,44, No.\,7}%               Set in Eng. 


%\str{\pageref{begin}--\pageref{end}}%  Set in Eng. 


\udk{514.753 } 


 


\author{\it Н.Е.\,Марюкова} 


% NB:   ^^ remove in Eng. 


\title{Об одной геометрической интерпретации\\


уравнения $\sin$-Гордона} 


 


\date{} 


% \pagestyle{myheadings}%         Set in Eng. 


\pagestyle{plain} %              Remove in Eng. 


\begin{document} 


% \thispagestyle{plain}%          Set in Eng. 


\maketitle 


\label{begin} 


 


{\bf 1. Введение.} В 1900 г. Д.\,Гильберт 


([1], приложение V) показал, что для угла $\psi(x,y)$ между 


асимптотическими 


линиями на поверхности с постоянной отрицательной кривизной $K_0$ в 


евклидовом пространстве выполняется уравнение 


\begin{equation} 


\psi_{x y}=-K_0\sin\psi. 


\end{equation} 


Это уравнение вытекает из более общего соотношения 


$\psi_{xy}+K\sin\psi=0$, 


полученного П.Л.\,Че\-бы\-ше\-вым в [2] для 


поверхности с произвольной кривизной $K(x,y)$ для угла $\psi(x,y)$ 


между линиями специальной сети на поверхности, которую теперь 


называют чебышевской сетью поверхности. Эта сеть характеризуется 


тем, что в каждом четырехугольнике, образованном линиями 


сети, противоположные стороны имеют одинаковую длину. Ясно, что 


нити куска ткани, натянутой на поверхность, образуют на ней 


чебышевскую сеть. Линии чебышевской сети принимаются за 


координатные линии на поверхности, за координаты $x$, $y$ принимаются 


длины дуг линий этой сети. Такие координаты на поверхности 


называются чебышевскими координатами. В чебышевских координатах 


первая квадратичная форма поверхности имеет вид 


\begin{equation} 


I=dx^2+2\cos\psi(x,y)dx\,dy+dy^2. 


\end{equation} 


Д.\,Гильберт показал, что чебышевская сеть на поверхности с постоянной 


отрицательной кривизной образована семействами асимптотических 


линий поверхности. 


 


Уравнение (1) называется уравнением sin-Гордона по аналогии с 


уравнением Клейна--Гордона 


$$ 


\psi_{u u}-\psi_{v v}=m^2\psi 


$$ 


для волновой функции $\psi(u,v)$ свободной частицы с массой $m$ и нулевым 


спином, если движение частицы характеризуется одной 


пространственной координатой $v$. Переменная $u$ в уравнении 


Клейна--Гордона пропорциональна времени. Уравнение (1) получается из 


уравнения Клейна--Гордона заменой функции $\psi$ в правой части на 


$\sin\psi$ 


и заменой переменных 


\begin{equation} 


u=x+y,\quad v=x-y. 


\end{equation} 


 


Результат Д.\,Гильберта получил дальнейшее развитие и 


обобщение в [3]--[6]. В 


этих работах связь метрики (2) и уравнения (1) была обобщена на случай 


произвольной метрики на поверхности с постоянной кривизной $K_0=-1$ и 


были получены дифференциальные уравнения, соответствующие 


известным метрикам на этой поверхности. 


 


В данной работе результат Д.\,Гильберта обобщается на более 


широкий класс пространств, вмещающих поверхность. Так как 


евклидово пространство является римановым пространством с 


кривизной, равной нулю, то возникает задача обобщения результата 


Д.\,Гильберта на случай риманова пространства с произвольной постоянной 


кривизной. Эти пространства рассмотрены в п.\,2. В п.\,3 работы 


доказано, что для угла между асимптотическими линиями поверхности с 


постоянной отрицательной гауссовой кривизной, по модулю не равной 


кривизне пространства, выполняется уравнение sin-Гордона. Если 


кривизна поверхности постоянна, отрицательна и по модулю равна 


кривизне пространства, то для угла между асимптотическими линиями 


поверхности выполняется уравнение 


\begin{equation} 


\psi_{v v}-\psi_{u u}=0. 


\end{equation} 


В п.\,4 рассмотрены чебышевские координаты на поверхности, в п.\,5 


приведен пример поверхности, для угла между асимптотическими 


линиями которой выполняется уравнение (4). 


\medskip 


 


{\bf 2. Римановы пространства постоянной кривизны.} Риманово 


пространство постоянной кривизны, равной $\frac{1}{r^2}$ или 


$-\frac{1}{r^2}$, является 


соответственно эллиптическим пространством, которое будем 


обозначать $S_3$, или гиперболическим пространством (пространством 


Лобачевского), которое будем обозначать $H_3$. Будем рассматривать 


пространства $S_3$ и $H_3$ как проективное пространство $P_3$ с абсолютом 


$$ 


(x^0)^2+\varepsilon\sum(x^i)^2=0, 


$$ 


где $\varepsilon=1$ для пространства $S_3$ и $\varepsilon=-1$ для 


пространства $H_3$. Тогда компоненты 


метрических тензоров этих пространств примут вид


\begin{equation} 


g_{00}=1,\quad g_{0i}=0,\quad 


g_{i j}=0\ \;\text{при}\ \; i\ne j,\quad 


g_{i i}=\varepsilon. 


\end{equation} 


Здесь и всюду в дальнейшем индексы пробегают значения


$i,j,k=1,2,3$; $I,J,K=0,1,2,3$. 


 


Продифференцируем равенства $e_I e_J=g_{I J}$ с учетом формул 


подвижного репера проективного пространства $d e_I=\omega^J_I e_J$, в 


результате 


чего получим соотношения $\omega^K_I g_{K J}+\omega^K_J g_{K I}=0$.


Учитывая (5), из этих соотношений получаем условия, 


которым удовлетворяют формы $\omega^J_I$,


\begin{equation} 


\omega^0_0=0,\quad \omega^0_i+\varepsilon\omega^i_0=0, 


\quad \omega^i_j+\omega^j_i=0. 


\end{equation} 


Из структурных уравнений проективного пространства 


$D\omega^I_J=\omega^K_J\wedge\omega^I_K$ с 


учетом соотношений (6) получаем структурные уравнения пространств 


$S_3$ и $H_3$. Полагая кривизну этих пространств равной 


$\varepsilon=\pm1$, находим 


формы кривизны этих пространств в виде 


$$ 


\Omega^j_i=D\omega^j_i-\omega^k_i\wedge\omega^j_k= 


-\varepsilon\omega^i_0\wedge\omega^j_0. 


$$ 


Обозначим 


\begin{equation} 


\omega^i_0=\omega^i. 


\end{equation} 


 


Пусть поверхность $V$ погружена в $S_3$ или $H_3$. В каждой 


точке $x=e_0\in V$ рассмотрим такой подвижный репер, что $e_1$, $e_2$ 


лежат в 


касательной плоскости к поверхности в точке $x\in V$. Тогда 


\begin{equation} 


\omega^3=0, 


\end{equation} 


откуда обычным образом получаем 


\begin{equation} 


\omega^3_1=a\omega^1+b\omega^2,\quad 


\omega^3_2=b\omega^1+c\omega^2. 


\end{equation} 


Поместим $e_1$ и $e_2$ на главных направлениях поверхности $V$ в точке 


$x$, 


тогда 


\begin{equation} 


b=0 


\end{equation} 


и гауссова кривизна поверхности, равная произведению главных 


кривизн, равна $K=a c$. 


\medskip 


 


{\bf 3. Поверхность с постоянной отрицательной гауссовой 


кривизной.} Пусть гауссова кривизна поверхности постоянна и 


отрицательна во всех точках этой поверхности, т.\,е. 


\begin{equation} 


K=-C^2,\quad C=\const\ne0. 


\end{equation} 


В этом случае уравнение асимптотических линий поверхности 


\begin{equation} 


a(\omega^1)^2+c(\omega^2)^2=0 


\end{equation} 


имеет два действительных различных корня. Обозначим через 


$\psi=2\varphi$


угол между асимптотическими линиями на поверхности, тогда 


\begin{equation} 


a=-C\tg\varphi,\quad c=C\ctg\varphi. 


\end{equation} 


 


Перейдем к локальным координатам $u$, $v$ в точке $x\in V$, полагая 


\begin{equation} 


\omega^1=A\,du,\quad \omega^2=B\,d v, 


\end{equation} 


где $A$, $B$ зависят от $u$, $v$. Дифференцируя соотношения (14) внешним 


образом с учетом структурных уравнений пространств $S_3$ и $H_3$ и формул 


(6), (7), (8) и (14), получаем 


\begin{equation} 


\omega^2_1=-\frac{A_v}{B}d u+\frac{B_u}{A}d v. 


\end{equation} 


Поступая аналогичным образом с уравнениями (9) и подставляя в них 


значения (10) и (13), получаем 


$$ 


-A_v\ctg\varphi=(A\tg\varphi)_v,\quad 


-B_u\tg\varphi=(B\ctg\varphi)_u. 


$$ 


Из полученных соотношений вытекает 


$$ 


\bigg(\ln\frac{A}{\cos\varphi} \bigg)_v=0,\quad 


\bigg(\ln\frac{B}{\sin\varphi} \bigg)_u=0 


$$ 


и, следовательно, можно положить 


\begin{equation} 


A=\cos\varphi,\quad B=\sin\varphi. 


\end{equation} 


Тогда (15) принимает вид


\begin{equation} 


\omega^2_1=\varphi_v du+\varphi_u dv. 


\end{equation} 


Дифференцируя равенство (17) внешним образом и находя внешний 


дифференциал формы $\omega^2_1$ из структурных уравнений пространств 


$S_3$ и 


$H_3$ с учетом формул (6)--(10) и (13)--(16), получаем


\begin{equation} 


\varphi_{v v}-\varphi_{u u}=(\varepsilon-C^2)\sin\varphi 


\cos\varphi. 


\end{equation} 


 


Поверхность в римановом пространстве имеет кривизны двух типов: 


внутреннюю кривизну $K_0$, равную секционной кривизне поверхности, 


рассматриваемой как двумерное риманово многообразие, и внешнюю,


или гауссову, кривизну $K$, равную произведению главных кривизн 


поверхности. Эти две кривизны связаны соотношением 


$$ 


K_0=K+\varepsilon, 


$$ 


где $\varepsilon$ --- кривизна риманова пространства. Заметим, что в 


евклидовом 


пространстве внутренняя и внешняя кривизны поверхности совпадают. В 


силу (11) выражение $\varepsilon-C^2$ в правой части уравнения (18) 


равно внутренней кривизне поверхности $\varepsilon-C^2=K_0$. Поэтому 


(18) можно переписать в виде 


\begin{equation} 


\varphi_{u u}-\varphi_{v v}=-K_0\sin\varphi\cos\varphi. 


\end{equation} 


Заменой $2\varphi=\psi$ и заменой координат (3) уравнение (19) при 


условии 


$K_0\ne0$ сводится к уравнению sin-Гордона (1), в котором величина $K_0$ 


означает внутреннюю кривизну поверхности, а при условии $K_0=0$ 


уравнение (19) сводится к уравнению (4) заменой $2\varphi=\psi$. 


 


Внутренняя кривизна поверхности может быть равна нулю только 


в эллиптическом пространстве $S_3$, т.\,к. в этом случае 


$C^2=\varepsilon=1$. В 


этом случае гауссова кривизна поверхности по модулю равна кривизне 


пространства. 


\medskip 


 


{\bf 4. Чебышевские координаты на поверхности.} Перейдем к 


координатам $x$, $y$ на поверхности с помощью замены координат (3). Из 


формул (14), (16) и (3) находим формы 


\begin{equation} 


\omega^1=\cos\varphi(x,y)(dx+dy),\quad 


\omega^2=\sin\varphi(x,y)(dx-dy). 


\end{equation} 


Тогда дифференциальное уравнение асимптотических линий (12) 


запишется в виде 


$$ 


2C\sin\psi(x,y)dx\,dy=0, 


$$ 


где $\psi(x,y)=2\varphi(x,y)$ --- угол между асимптотическими линиями. Из 


этого 


уравнения вытекает, что координатные линии $x=\const$ и $y=\const$ на 


поверхности являются асимптотическими линиями. Подставляя в 


первую квадратичную форму $I=(\omega^1)^2+(\omega^2)^2$ значения (20), 


находим, что в 


координатах $x$, $y$ первая квадратичная форма поверхности имеет вид (2), 


координаты $x$, $y$ означают длины дуг координатных линий. Таким 


образом, координаты $x$, $y$ являются чебышевскими координатами на 


поверхности. В этих координатах уравнение (19) при $\psi=2\varphi$ имеет 


вид 


(1), а уравнение (4) --- вид 


\begin{equation} 


\psi_{x y}=0. 


\end{equation} 


Таким образом, доказана следующая 


 


\begin{thmnonum} 


В трехмерном римановом пространстве постоянной 


кривизны дана поверхность постоянной отрицательной гауссовой 


кривизны $K$, отнесенная к чебышевским координатам. Если гауссова 


кривизна поверхности по модулю не равна кривизне пространства, то 


для угла $\psi$ между асимптотическими линиями поверхности 


выполняется уравнение $\sin$-Гордона $(1)$, в котором величина $K_0$ 


означает 


внутреннюю кривизну поверхности. Если гауссова кривизна поверхности 


по модулю равна кривизне пространства, то для угла между 


асимптотическими линиями выполняется уравнение $(21)$. 


\end{thmnonum} 


 


{\bf 5. Пример поверхности с нулевой внутренней кривизной.} 


Приведем пример поверхности в эллиптическом пространстве $S_3$, 


внутренняя кривизна которой в каждой точке равна нулю. Это так 


называемая поверхность Клиффорда ([7], с.\,93--95). Ее можно определить 


как множество точек, равноотстоящих от данной прямой, называемой 


осью, и как множество прямых, паратактичных оси, т.\,е. находящихся от 


оси на постоянном расстоянии. Так как через каждую точку 


пространства $S_3$ проходят две прямые, паратактичные данной прямой, 


то поверхность Клиффорда имеет два семейства прямолинейных 


образующих, являющихся асимптотическими линиями этой 


поверхности. Так как угол между образующими, проходящими через одну 


точку, находящуюся на расстоянии $\delta$ от оси, равен $2\delta$, а 


расстояние 


$\delta$ постоянно для поверхности Клиффорда, то угол $\psi=2\delta$ 


между 


асимптотическими линиями постоянен и, следовательно, удовлетворяет 


уравнению (21). 


 


В ([7], с.\,94--95) показано, что поверхность Клиффорда изометрична 


евклидовой плоскости и при развертывании на плоскость эта 


поверхность изображается в виде ромба с острым углом $\psi=2\delta$ и 


длиной стороны $\pi$. 


 


\begin{thebibliography}{9} 


 


\bibitem{1} 


Гильберт Д. {\em Основания геометрии}. -- М.--Л.: ОГИЗ, 1948. -- 491 с. 


\bibitem{2} 


Чебышев П.Л. {\em О кройке одежды} // ПСС. -- М.--Л.: АН СССР, 1951. -- 


Т.\,5. -- С.\,165--170. 


\bibitem{3} 


Позняк Э.Г. {\em Геометрические исследования, связанные с уравнением 


$z_{x{}y}=\sin{}z$} // 


Итоги науки и техн. ВИНИТИ. Пробл. геометрии. -- 1977. -- Т.\,8. 


-- С.\,225--241. 


\bibitem{4} 


Позняк Э.Г., Попов А.Г. {\em Геометрия уравнения $\sin$-Гордона} // Итоги 


науки и техн. ВИНИТИ. Пробл. геометрии. -- 1991. -- Т.\,23. -- 


С.\,99--130. 


\bibitem{5} 


Позняк Э.Г., Попов А.Г. {\em Геометрия Лобачевского и физика} // Изв. 


вузов. Математика. -- 1994. -- \No\,3. -- С.\,44--49. 


\bibitem{6} 


Попов А.Г. {\em Геометрический метод точного интегрирования 


эллиптического уравнения Лиувилля $\Delta{}u=e^u$} // Вестн. Моск. ун-та. 


Сер. матем. механ. -- 1995. -- \No\,3. -- С.\,82--84. 


\bibitem{7} 


Розенфельд Б.А. {\em Неевклидовы пространства}. -- М.: Наука, 1969. -- 


547 с. 


 


\end{thebibliography} 


 


\vskip 0.5cm 


 


\post{\it Брянский государственный}{\it Поступила} 


\post{\it педагогический университет}{20.07.1998} 


 


\label{end} 


\end{document} 


